\section{Einleitung}
\label{sec:introduction}

\subsection{Motivation und Problemstellung}

Neuronale Netze werden zunehmend auf ressourcenbeschraenkter Hardware ausgeführt, etwa auf eingebetteten Systemen, mobilen Geräten und Edge-Beschleunigern. Auf solcher Hardware sind die Inferenzkosten eines vollpräzisen Modells oft nicht mehr tragbar. Die Quantisierung reduziert diese Kosten, indem die Gewichte und Aktivierungen eines Netzes von 32-Bit-Gleitkommazahlen (FP32) auf niedriger aufgelöste, diskrete Werte abgebildet werden, typischerweise auf 8 Bit oder weniger. 
\par\medskip

Eine besonders hardware-attraktive Variante ist die Power-of-Two-Quantisierung (PoT). Hierbei werden die Gewichte nicht auf ein gleichmäßiges, uniformes Gitter abgebildet, sondern auf Zweierpotenzen. Der Vorteil liegt in der Arithmetik, da eine Multiplikation mit einer Zweierpotenz einem einfachen Bitshift entspricht. Dadurch können die sonst dominierenden Multiplizierer-Schaltkreise in der Hardware entfallen \citep{miyashita_pot_2016, przewlocka_rus_pot_2022}. Dieser Effizienzgewinn ist jedoch mit einem Nachteil verbunden. Da das PoT-Gitter logarithmisch und nicht linear ist, variiert der Quantisierungsfehler stark und unvorhersehbar mit den tatsächlichen Gewichtswerten einer Schicht. 
\par\medskip

Um zu entscheiden, welche Schichten eines Netzes eine niedrigere Bit-Breite vertragen und welche nicht, hat sich die Verwendung von Krümmungsinformation aus der zweiten Ableitung des Verlustes etabliert, also der Hessian-Matrix. Ein Ansatzes ist die HAWQ-Familie \citep{dong_hawq_2019, dong_hawqv2_2020}. HAWQ rangiert die Schichten nach ihrem größten Hessian-Eigenwert, während HAWQ-V2 dies zu einem Produkt aus Hessian-Spur und quadrierter Stoerungsgroeße, $\mathrm{Tr}(H)\cdot\|\Delta W\|^2$, verfeinert, motiviert durch eine Taylor-Entwicklung des Verlustes zweiter Ordnung. Beide Arbeiten nutzen diese Formulierung, um eine Bit-Zuteilung über das gesamte Netz zu optimieren. Validiert wird dabei die End-to-End-Genauigkeit nach der Zuteilung und nicht die Vorhersagegüte des Scores auf Ebene einzelner Schichten.
\par\medskip

\subsection{Forschungsfragen}

Das Projekt verfolgt drei Forschungsfragen:

\begin{enumerate}
\item Wie prägen PTQ und QAT die schichtweise Verlust-Krümmung (gemessen über die Hessian-Spur) von CNNs unter PoT-Quantisierung jeweils unterschiedlich?
\item Sagt die schichtweise Verlust-Krümmung (Hessian-Spur) vorher, welche Schichten tatsächlich den größten PoT-Gewichts-Quantisierungsschaden erleiden, und wie haengen diese beiden Auswahlen (vorhergesagt vs.\ tatsächlich) von Architektur und Datensatz ab?
\item In welchem Ausmaß ist der PoT-induzierte Genauigkeitsverlust auf Gewichts- gegenüber Aktivierungsquantisierung zurückführbar, und für welches Regime ist demzufolge die Hessian-basierte Sensitivitätsanalyse anwendbar?
\end{enumerate}

Die zweite Forschungsfrage ist die zentrale Fragestellung und der Kern des daraus entstandenen Workshop-Papers (\S\ref{sec:conclusion}). Die restlichen beiden sind breiter angelegt und werden in diesem Bericht vollständig behandelt, obwohl sie im Paper aus Platzgründen nur gestreift oder ausgeklammert werden.